\documentclass[a4paper,12pt]{article}
\usepackage[spanish,activeacute]{babel}
\usepackage[latin1]{inputenc}

%% Para las columnas
\usepackage{multicol}
\usepackage{amssymb,amsmath,eepic,fancyhdr}
\usepackage{latexsym}
\usepackage{datenumber}


%% Para numeraci?n autom?tica
\newcounter{nroproblema}
\setcounter{nroproblema}{1}
\newcounter{nroapartado}
\setcounter{nroapartado}{1}

\newcommand{\n}{%
    \vspace{.1in}\noindent{\bf \arabic{nroproblema}.\ }%
    \addtocounter{nroproblema}{1}%
    \setcounter{nroapartado}{1}%
    }%

\newcommand{\nn}{%
    \vspace{.1in}\noindent \hskip.4cm {\bf \arabic{nroproblema}.\ }%
    \addtocounter{nroproblema}{1}%
    \setcounter{nroapartado}{1}%
    }%


\newcommand{\ap}{\vspace*{0.2cm}\hspace*{0.2cm}%
   {\bf \alph{nroapartado})}
   \addtocounter{nroapartado}{1}%
   }

 
%marcos by Marco
\newcommand{\pd}[1]{\frac{\partial}{\partial #1}} %\pd{x}
\newcommand{\NN}{\ensuremath{\mathbb N}}
\newcommand{\ZZ}{\ensuremath{\mathbb Z}}
\newcommand{\CC}{\ensuremath{\mathbb C}}
\newcommand{\RR}{\ensuremath{\mathbb R}}
\newcommand{\TT}{\ensuremath{\mathbb T}}
\newcommand{\g}{\ensuremath{\frak{g}}}
\newcommand{\h}{\ensuremath{\frak{h}}}
\newcommand{\ft}{\ensuremath{\frak{t}}}
\newcommand{\vX}{\frak{X}}
\newcommand{\cB}{\mathcal{B}}
\newcommand{\cN}{\mathcal{N}}
\newcommand{\cL}{\mathcal{L}}
\newcommand{\cD}{\mathcal{D}} 
 

%% m�rgenes

% Anterior:

%\setlength{\unitlength}{1mm} \addtolength{\voffset}{-15mm}
%\addtolength{\textheight}{30mm} \addtolength{\hoffset}{-20mm}
%\addtolength{\textwidth}{45mm}

\setlength{\unitlength}{1mm} \addtolength{\voffset}{-22mm}
\addtolength{\textheight}{40mm} \addtolength{\hoffset}{-22mm}
\addtolength{\textwidth}{46mm}




\begin{document}

\noindent
{\sffamily\large Marco Zambon 
\hfill  Master UAM, 2013-14  \\Geometr\'ia simpl\'ectica y de Poisson}


 

\begin{center}\bfseries\large\textsf{Hoja 5    \vspace{-2mm}}
\end{center}
 
 
\n
\ap Sea $G$ un grupo de Lie. Sea $(M,\omega)$ una variedad simpl\'ectica y $\Phi$ una accion hamiltoniana de $G$ en $(M,\omega)$  con aplicaci\'on momento $J\colon M\to \g^*$.

Sea $H$ un subgrupo de Lie de $G$. A partir de $J$, construye una aplicaci\'on momento para la acci\'on $\Phi|_H$ de $H$ en $(M,\omega)$.


\ap
Sea $G$ un grupo de Lie. Sean $(M_i,\omega_i)$ variedades simpl\'ecticas y $\Phi_i$ acciones hamiltonianas de $G$ en $(M_i,\omega_i)$  con aplicaci\'on momento $J_i$ ($i=1,2$).

Demuestra que la acci\'on diagonal de $G$ en $(M_1\times M_2,\pi_1^*\omega_1+\pi_2^*\omega_2)$ (es decir, $g\mapsto \Phi_1(g) \times \Phi_2(g)$) es hamiltoniana, y describe una aplicaci\'on momento. Aqu\'i $\pi_i\colon M_1\times M_2\to M_i$ es la proyecci\'on.
 
  \noindent\textbf{Consejo:} Considera primero la acci\'on $\Phi_1\times \Phi_2$ de $G\times G$ en $M_1\times M_2$, despu\'es utiliza a) y la aplicaci\'on diagonal $\Delta\colon \g\to \g\times \g, v\mapsto (v,v)$.
  
  


 
 \n  Sea $(W,\omega)$ un espacio vectorial simpl\'ectico, y $Sp(W,\omega)$ su grupo de simplectomorf\'ismos. Sea $G$ un subgrupo de Lie de 
$Sp(W,\omega)$. 
Comprueba que una aplicaci\'on momento para la acci\'on can\'onica de  $G$ en $(W,\omega)$ es   $J\colon W\to \g^*$, cuyas componentes son dadas como sigue: para todo $A\in \g$,  $$J^A\colon W\to \RR, w\mapsto-\frac{1}{2}\omega(w,Aw).$$
  
  \noindent\textbf{Nota:} puedes mirar la p\'agina 66 de los apuntes de Meinrenken, donde Meinrenken no demuestra la equivarianza.
  
  
  \n   Sean $\lambda_1,\dots,\lambda_n\in \ZZ$. Considera la acci\'on del circulo $U(1)$ en $\CC^n$ por  $$e^{i\theta}\cdot (z_1,\dots,z_n)= (e^{i\theta\lambda_1}z_1,\dots,e^{i\theta\lambda_n}z_n).$$
Encuentra la \'unica aplicaci\'on momento para esta acci\'on que lleva $0\in \CC^n$ al origen del dual del \'algebra  de Lie, de dos maneras distintas:

\ap    utilizando el problema 1b)

\ap utilizando el problema 2.
  
  
\n  \ap Considera una acci\'on de un grupo de Lie $G$   sobre una variedad 
 simpl\'ectica  \emph{compacta} $(M,\omega)$. Demuestra que si existe $v\in \g$ tal que  $(v_M)|_p\neq 0$ para todo  $p\in M$, entonces no existe ninguna aplicaci\'on momento para esta acci\'on.
 
% \ap Deduce que toda acci\'on sin puntos fijos de $U(1)$ en una variedad 
 %simpl\'ectica  \emph{compacta}   no tiene aplicaci\'on momento.
%Me: need to know that d/dt exp(tv)p=v_M...for all t 
 
  \ap Demuestra que toda  acci\'on de  $G=SL_2(\RR)$ sobre una variedad 
 simpl\'ectica   $(M,\omega)$ por simplectomorfismos admite una \'unica aplicaci\'on momento. 
  
  \noindent\textbf{Consejo:} usa un problema de la hoja 4.
  
\n [3 puntos]
Sea $(M,\omega)$ una variedad simpl\'ectica y considera una acci\'on hamiltoniana de un grupo de Lie $G$ en $(M,\omega)$  con aplicaci\'on momento $J$. Demuestra que la equivarianza de $J$ implica que
$$\cL_{v_M}J^w=J^{[v,w]}$$
para todo $v,w\in \g$. Aqu\'i $J^v\colon M\to \RR$ es la componente de $J$ dada por $v$.


  
   \end{document}
  \noindent\textbf{Consejo:}
\noindent\textbf{Nota:}